\documentclass[11pt]{article}
\usepackage[T1]{fontenc}
\usepackage{lmodern}
\usepackage[margin=1in]{geometry}
\usepackage{amsmath,amssymb,amsthm,mathtools}
\usepackage{microtype}
\usepackage{enumitem}
\usepackage[colorlinks=true,linkcolor=blue,citecolor=blue,urlcolor=blue]{hyperref}
\usepackage{bookmark}

\newcommand{\Pbb}{\mathbb P}
\newcommand{\Ebb}{\mathbb E}
\newcommand{\R}{\mathbb R}
\newcommand{\Z}{\mathbb Z}
\newcommand{\T}{\mathbb T^2}
\newcommand{\cM}{\mathcal M}
\newcommand{\supp}{\operatorname{supp}}
\newcommand{\diver}{\operatorname{div}}
\newcommand{\Var}{\operatorname{Var}}
\newcommand{\Lip}{\operatorname{Lip}}
\newcommand{\dd}{\,\mathrm d}
\newcommand{\norm}[1]{\left\lVert #1\right\rVert}
\newcommand{\abs}[1]{\left\lvert #1\right\rvert}
\newcommand{\eps}{\varepsilon}
\DeclareMathOperator*{\esssup}{ess\,sup}
\newtheorem{theorem}{Theorem}[section]
\newtheorem{lemma}[theorem]{Lemma}
\newtheorem{proposition}[theorem]{Proposition}
\newtheorem{corollary}[theorem]{Corollary}
\theoremstyle{remark}
\newtheorem{remark}[theorem]{Remark}
\numberwithin{equation}{section}
\setlist[enumerate]{itemsep=3pt,topsep=5pt}

\title{Simultaneous quadratic-cost and bottleneck bounds\\for planar random matching}
\author{}
\date{\today}

\begin{document}
\maketitle

\begin{abstract}
For two independent samples of \(N\) uniform points in the unit square,
we construct a single matching whose total squared distance is
\(O(\log N)\) and whose longest edge is
\(O((\log N)^{3/4}/\sqrt N)\), with failure probability
\(O((\log N)^{-2})\). The matching can be chosen measurably so that its
expected quadratic cost is \(O(\log N)\) and its bottleneck bound fails
with probability at most \(C\exp(-c(\log N)^{3/2})\). We also obtain
these conclusions for a random sample matched to fixed evenly spread
points.

The proof combines the classical bottleneck matching theorem with a
Fourier estimate for smoothed empirical measures. Lipschitz truncation
of a scalar stream function gives one divergence solution with both
an \(L^2\) bound and an \(L^\infty\) bound. A continuity equation with
uniformly positive densities converts this field into a coupling that
retains both estimates. Rounding preserves the distance constraint and
produces a permutation with the required quadratic cost.
\end{abstract}

\section{Statement of the main result}\label{sec:statement}

Let \(X_1,\ldots,X_N,Y_1,\ldots,Y_N\) be independent and uniformly distributed on \([0,1]^2\), and write \(d(x,y)=|x-y|\) for Euclidean distance. The set of permutations of \(\{1,\ldots,N\}\) is denoted by \(\mathfrak S_N\). We prove the following statement.

\begin{theorem}\label{thm:main}
    There are universal constants \(C_0,C_1,C_2>0\) such that, for all sufficiently large \(N\), with probability at least \(1-C_2/(\log N)^2\), there exists \(\pi\in\mathfrak S_N\) satisfying
    \[
        \sum_{i=1}^N d(X_i,Y_{\pi(i)})^2\le C_0\log N, \qquad \max_{1\le i\le N}d(X_i,Y_{\pi(i)}) \le \frac{C_1(\log N)^{3/4}}{\sqrt N}.
    \]
\end{theorem}
\begin{remark}
For the same permutation, Cauchy--Schwarz gives
\[
 \sum_{i=1}^N d(X_i,Y_{\pi(i)})
 \le \sqrt{N\sum_{i=1}^N d(X_i,Y_{\pi(i)})^2}
 \le \sqrt{C_0}\sqrt{N\log N}.
\]
Thus, with \(L=\max\{\sqrt{C_0},C_1\}\),
\[
 \lim_{N\to\infty}\Pbb\left(
 \exists\pi\in\mathfrak S_N:\quad
 \begin{aligned}
  \sum_{i=1}^N d(X_i,Y_{\pi(i)})
     &\le L\sqrt{N\log N},\\
  \max_{1\le i\le N}d(X_i,Y_{\pi(i)})
     &\le \frac{L(\log N)^{3/4}}{\sqrt N}
 \end{aligned}
 \right)=1.
\]
This is the simultaneous conclusion in \cite[Problem~5.3]{Tal26}.
The corresponding fixed-target formulation in
\cite[Question~4.9.1]{Tal21} is treated in
Corollary~\ref{cor:grid}, including its expected total-distance bound.
The full ultimate matching conjecture
\cite[Research Problem~17.1.2]{Tal21} asks for stronger coordinatewise
exponential integrability. Even its symmetric exponent-four case is
not established by the quadratic-cost and maximal-distance estimates
proved here.
\end{remark}

The total-distance and bottleneck orders separately are classical;
see \cite[Theorems~4.5.1 and~4.7.1]{Tal21}. The unconstrained
quadratic matching cost is also understood: Ambrosio, Stra, and
Trevisan \cite[Theorem~1.1]{AST19} prove that its expectation is
asymptotic to \((2\pi)^{-1}\log N\) for two uniform square samples.
The point of Theorem~\ref{thm:main} is to obtain the quadratic order
while imposing the bottleneck constraint on the same permutation.

Ouimet's preprint \cite[Theorem~2.1]{Oui26} gives simultaneous
total-distance and bottleneck bounds. Its planar flow interpolation
\cite[Lemmas~2.9--2.10]{Oui26} uses a closely related stream-function
truncation. The present quadratic estimate is stronger than that
theorem's stated conclusion: multiplying its total-distance and
maximal-distance bounds gives only
\[
 \sum_i d(X_i,Y_{\pi(i)})^2
 \le \bigl(\max_i d(X_i,Y_{\pi(i)})\bigr)
          \sum_i d(X_i,Y_{\pi(i)})
 =O((\log N)^{5/4}).
\]
Here the conversion through smooth densities retains the
\(L^2\) energy all the way to the final matching.
We state the deterministic compatibility result in
Proposition~\ref{prop:continuous}.
In \S\ref{sec:further}, we also derive expected-cost bounds and the
fixed-target version from the same estimates.


\section{Reduction to the torus and proof overview}
\label{sec:proof-structure}

Let \(\T=(\R/2\Z)^2\) be the flat torus with geodesic distance
\(d_{\T}\). Its ordinary Lebesgue measure has total mass \(4\), so
the uniform probability density is \(1/4\). We first state the torus
version of the matching theorem.

\begin{theorem}[Matching on the torus]
\label{thm:torus}
Let \(X_1,\ldots,X_N,Y_1,\ldots,Y_N\) be independent uniform
points on \(\T\). There are universal constants \(C_0,C_1,C_2>0\)
such that, for all sufficiently large \(N\), with probability at least
\(1-C_2/(\log N)^2\), there exists \(\pi\in\mathfrak S_N\)
satisfying
\[
 \sum_{i=1}^N d_{\T}(X_i,Y_{\pi(i)})^2\le C_0\log N,
 \qquad
 \max_{1\le i\le N}d_{\T}(X_i,Y_{\pi(i)})
 \le\frac{C_1(\log N)^{3/4}}{\sqrt N}.
\]
\end{theorem}

\begin{proof}[Proof of Theorem~\ref{thm:main} assuming Theorem~\ref{thm:torus}]
Let \(X_1,\ldots,X_N,Y_1,\ldots,Y_N\) have the torus law in
Theorem~\ref{thm:torus}. We fold these samples onto the square using
the map \(F:\T\to[0,1]^2\) defined by
\[
 F(z_1,z_2)=
 \bigl(d_{\R/2\Z}(z_1,0),\,d_{\R/2\Z}(z_2,0)\bigr).
\]
On the representatives \([-1,1)^2\), this is simply
\(F(z_1,z_2)=(|z_1|,|z_2|)\). A uniform torus point has independent
coordinates uniform on \([-1,1)\), whose absolute values are
independent and uniform on \([0,1]\). Consequently,
\(F(X_1),\ldots,F(X_N),F(Y_1),\ldots,F(Y_N)\) have exactly the
joint distribution of the square samples in Theorem~\ref{thm:main}.

Distance to a fixed point is a \(1\)-Lipschitz function on any metric
space. Applying this fact to each circle coordinate gives
\[
 |F(z)-F(z')|^2
 \le\sum_{\ell=1}^2d_{\R/2\Z}(z_\ell,z'_\ell)^2
 =d_{\T}(z,z')^2.
\]
Thus every matching of the torus samples induces a matching of the
folded samples with no increase in any edge length. On the event in
Theorem~\ref{thm:torus}, the same permutation \(\pi\) therefore satisfies
\[
 \sum_{i=1}^N|F(X_i)-F(Y_{\pi(i)})|^2\le C_0\log N,
 \qquad
 \max_{1\le i\le N}|F(X_i)-F(Y_{\pi(i)})|
 \le\frac{C_1(\log N)^{3/4}}{\sqrt N}.
\]
The event that the folded samples admit such a permutation contains
the torus matching event, so its probability is at least
\(1-C_2/(\log N)^2\). The equality of the joint sample laws proves
Theorem~\ref{thm:main} with the same constants.
\end{proof}

It remains to prove Theorem~\ref{thm:torus}. We sketch its proof,
emphasizing how the quadratic-cost and maximal-distance bounds are
obtained for a single matching. Throughout the rest of the paper,
\(X_1,\ldots,X_N,Y_1,\ldots,Y_N\) denote the independent uniform
torus samples in that theorem. We first construct a coupling of two
smooth densities and then transfer its bounds to a permutation.

We smooth each empirical measure at radius
\[
 h=A\sqrt{\frac{\log N}{N}},
\]
where \(A\) is a sufficiently large fixed constant. With probability
tending to one, the resulting densities \(\rho_0,\rho_1\) are bounded
above and below by positive constants. Two separate constructions then
give solutions of the equation \(\diver u=\rho_0-\rho_1\). A Fourier
calculation yields a field \(u\) with
\(\norm u_2^2=O(\log N/N)\). The bottleneck matching theorem, followed
by a common smoothing of the matched endpoints, gives a coupling with
distance cutoff \(r=O((\log N)^{3/4}/\sqrt N)\). From this coupling
we construct another field \(v\), with the same divergence, satisfying
\(\norm v_\infty=O(r)\).

The main deterministic step combines these two fields. After subtracting
their spatial means, the divergence-free difference \(u-v\) is the
rotated gradient of a scalar stream function. On the set where the
maximal function of \(\abs u\) is controlled, this stream function is
Lipschitz with constant \(O(\norm v_\infty)\). Extending its restriction
from that set gives a field \(w\) with the same divergence and with
\[
 \norm w_2^2\le C\norm u_2^2,
 \qquad \norm w_\infty\le C\norm v_\infty.
\]
The area of the modified set controls the additional energy. This is
the step that makes the two estimates compatible.

We next interpolate linearly between \(\rho_0\) and \(\rho_1\), using
\(w\) as a fixed flux. The lower density bound controls both the speed
and the quadratic action of the resulting flow. Its endpoint coupling
therefore has quadratic cost \(O(\log N/N)\) and distance cutoff
\(O(r)\). Passing from the smoothed measures back to the empirical
measures adds at most \(2h\) to the cutoff and \(O(h^2)\) to the
normalized quadratic cost. These errors have the required sizes because
\(h=o(r)\) and \(h^2=O(\log N/N)\). A decomposition into permutation
matrices then selects one permutation with both bounds.

The deterministic inputs are collected in \S\ref{sec:preliminaries},
culminating in Proposition~\ref{prop:continuous}.
In \S\ref{sec:random-inputs} we establish the density and energy
estimates and apply the Leighton--Shor bottleneck theorem
\cite{LS89}, in the form stated in \cite[Theorem~4.7.1]{Tal21}.
The rounding argument in \S\ref{sec:rounding} then completes the
proof of Theorem~\ref{thm:torus}.
The further consequences in \S\ref{sec:further} use a matching
that minimizes quadratic cost subject to a fixed distance cutoff.
The density, energy, and bottleneck estimates hold together with probability at least
\(1-C/(\log N)^2\), which gives the stated failure probability.
Constants denoted by \(C,c>0\) may change between occurrences.
They are independent of \(N\), and may depend on the fixed mollifier
introduced below.

\section{Preliminary inputs}
\label{sec:preliminaries}

This section develops the deterministic ingredients of the proof. The
first lemma combines two solutions of the same divergence equation.
The second constructs a bounded solution from a coupling with a
distance cutoff. The third turns a field satisfying both bounds into a
single coupling, leading to Proposition~\ref{prop:continuous}.
The common setting and normalizations are fixed here for the remainder
of the paper.

We use the torus and geodesic distance defined in
\S\ref{sec:proof-structure}. All spatial norms and integrals below
are taken with respect to ordinary Lebesgue measure on \(\T\), whose
total mass is \(4\); the uniform probability density is therefore
\(1/4\). For an integrable scalar function or
vector field \(g\), its spatial mean is
\(\overline g=\frac14\int_{\T}g\). All divergence identities are
interpreted distributionally. For integrable \(f\), the identity
\(\diver u=f\) means
\[
 -\int_{\T}u\cdot\nabla\varphi\,\mathrm dx
 =\int_{\T}f\varphi\,\mathrm dx
 \qquad\text{for every smooth periodic test function }\varphi.
\]
A coupling of two probability densities is a probability measure on
\(\T\times\T\) with those densities as its marginals. We say it has
distance cutoff \(r\) when its support is contained in
\(\{(x,y):d_{\T}(x,y)\le r\}\).

Fix a nonnegative function \(\eta\in C_c^\infty(\R^2)\), supported
in the unit ball and with integral one. For \(0<s<1/8\), write
\[
 \eta_s^{\R}(z)=s^{-2}\eta(z/s),\qquad
 \eta_s(z)=\sum_{m\in\Z^2}\eta_s^{\R}(z+2m).
\]
Thus \(\eta_s\) is a smooth probability kernel on \(\T\), supported
where \(d_{\T}(z,0)\le s\). Convolution on the torus uses ordinary
Lebesgue measure, with
\((\eta_s*g)(x)=\int_{\T}\eta_s(x-y)g(y)\,\mathrm dy\).
For a finite measure the same notation denotes convolution against that
measure. Integrals over \(\R^2\) use \(\eta_s^{\R}\), and periodic
functions appearing in such integrals are evaluated periodically.
No symmetry assumption on \(\eta\) is needed.

\subsection{Combining two divergence solutions}
\label{sec:interpolation}

The purpose of the next lemma is to retain an \(L^2\) bound from one
divergence solution and an \(L^\infty\) bound from another. We first
recall the analytic estimates used in its proof. For an integrable
function \(g\), define the periodic Hardy--Littlewood maximal function
by
\[
 \cM g(x)=\sup_{s>0}\frac1{|B_{\T}(x,s)|}
                  \int_{B_{\T}(x,s)}|g(y)|\,\mathrm dy.
\]
Lebesgue differentiation gives \(\abs g\le\cM g\) almost everywhere,
and \(\cM g\le\norm g_\infty\) for bounded \(g\). The weak maximal
inequality and \((\cM g)^2\le\cM(|g|^2)\) imply
\[
 \bigl|\{x\in\T:\cM g(x)>t\}\bigr|
 \le \frac{C}{t^2}\norm g_2^2,
 \qquad g\in L^2(\T),\quad t>0.
\]
We also use the pointwise Sobolev estimate
\begin{equation}
\label{eq:sobolev-pointwise}
 |\psi(x)-\psi(y)|
 \le C d_{\T}(x,y)
 \bigl(\cM(|\nabla\psi|)(x)+\cM(|\nabla\psi|)(y)\bigr),
 \qquad \psi\in H^1(\T),
\end{equation}
for every pair of Lebesgue points \(x,y\) of the chosen Sobolev
representative. In particular, a bound on the maximal
function of the gradient makes the restriction of \(\psi\) Lipschitz.

The Euclidean version and its use in Lipschitz approximation are given
in \cite[proof of Lemma~10, p.~236]{Haj00}. The same argument applies
on the fixed torus. To see this, let \(\psi_B\) denote the average of
\(\psi\) over a ball \(B\). Poincar\'e's inequality, followed by
telescoping the averages over \(B_{\T}(x,2^{-j}s)\), gives
\[
 |\psi(x)-\psi_{B_{\T}(x,s)}|
 \le Cs\,\cM(|\nabla\psi|)(x).
\]
For \(s=d_{\T}(x,y)\), compare the averages over
\(B_{\T}(x,s)\) and \(B_{\T}(y,s)\) with the average over
\(B_{\T}(x,3s)\). The volume-doubling property and Poincar\'e's
inequality bound their difference by the right-hand side of
\eqref{eq:sobolev-pointwise}. The preceding telescoping estimate at
both endpoints then proves that inequality. Large radii are covered by
the global Poincar\'e inequality on \(\T\). The periodic weak maximal
inequality follows from the usual covering argument on this compact
doubling space. All constants depend only on the fixed torus.

\begin{lemma}[Simultaneous bounds for a divergence solution]
\label{lem:interpolate}
Suppose \(u\in L^2(\T;\R^2)\) and
\(v\in L^\infty(\T;\R^2)\) have the same distributional divergence.
There exists \(w\in L^2(\T;\R^2)\cap L^\infty(\T;\R^2)\) such
that
\[
 \diver w=\diver u,\qquad
 \norm w_2^2\le C\norm u_2^2,\qquad
 \norm w_\infty\le C\norm v_\infty.
\]
\end{lemma}

\begin{proof}
We first remove the constant parts of the fields. Replacing \(u,v\)
by \(u-\overline u,v-\overline v\) preserves their divergences, and
\[
 \norm{u-\overline u}_2^2=\norm u_2^2-4|\overline u|^2
 \le\norm u_2^2,\qquad
 \norm{v-\overline v}_\infty\le2\norm v_\infty.
\]
It is therefore enough to prove the result for mean-zero fields, which
we continue to denote by \(u,v\). This reduction removes the constant
divergence-free fields that cannot be represented by a periodic stream
function.

The field \(u-v\) belongs to \(L^2\), has zero mean, and is
divergence free. Hence it is a rotated gradient:
\[
 u-v=\nabla^\perp\psi,
 \qquad \nabla^\perp\psi=(-\partial_2\psi,\partial_1\psi),
 \qquad \psi\in H^1(\T).
\]
For completeness, use the Fourier convention
\[
 \widehat g(k)=\frac14\int_{\T}g(z)e^{-i\pi k\cdot z}\,\mathrm dz,
 \qquad k\in\Z^2.
\]
For every \(k\ne0\), the divergence condition says that
\(\widehat{(u-v)}(k)\) is orthogonal to \(k\). With
\(k^\perp=(-k_2,k_1)\), define
\[
 \widehat\psi(k)=
 \frac{k^\perp\cdot\widehat{(u-v)}(k)}{i\pi|k|^2}
 \quad(k\ne0),\qquad \widehat\psi(0)=0.
\]
The Fourier coefficient of \(\nabla^\perp\psi\) is then exactly
\(\widehat{(u-v)}(k)\). Parseval's identity gives
\(\norm{\nabla\psi}_2=\norm{u-v}_2\), and the absence of the zero
mode also gives \(\psi\in L^2\). Thus this formula defines the
required real-valued, mean-zero function in \(H^1(\T)\).

Write \(B=\norm v_\infty\). If \(B=0\), the common divergence is
zero and we take \(w=0\). Suppose \(B>0\), and define
\[
 G=\{x\in\T:\cM(|u|)(x)\le B\},
\]
intersected with the full-measure set where the representatives and
\eqref{eq:sobolev-pointwise} are valid. On \(G\), the stream-function
identity implies
\[
 \cM(|\nabla\psi|)
 \le\cM(|u|)+\cM(|v|)\le2B.
\]
It follows from \eqref{eq:sobolev-pointwise} that \(\psi|_G\) is
\(CB\)-Lipschitz with respect to \(d_{\T}\).

If \(G\) is nonempty, extend this restriction by the McShane formula
\cite[Lemma~5]{Haj00}:
\[
 \widetilde\psi(x)
 =\inf_{y\in G}\{\psi(y)+CB\,d_{\T}(x,y)\}.
\]
Here the constant is chosen to dominate the Lipschitz constant of the
restriction. The infimum is finite: after fixing one point of \(G\),
the Lipschitz bound and the finite diameter of \(\T\) bound \(\psi\)
above and below on \(G\). The same bound shows that
\(\widetilde\psi=\psi\) on \(G\) and that
\(\widetilde\psi\) is globally \(CB\)-Lipschitz on \(\T\).
In particular, it is periodic and
\(\norm{\nabla\widetilde\psi}_\infty\le CB\). If \(G\) is empty,
set \(\widetilde\psi=0\).

Define \(w=v+\nabla^\perp\widetilde\psi\). Distributional derivatives
commute, so \(\diver\nabla^\perp\widetilde\psi=0\). Therefore
\(\diver w=\diver v=\diver u\), and
\(\norm w_\infty\le CB\).
Moreover, locality of weak derivatives gives
\(\nabla\widetilde\psi=\nabla\psi\) almost everywhere on \(G\),
and consequently \(w=u\) there. This locality property applies to a
measurable agreement set. Indeed, the difference of the two Sobolev
functions has an absolutely continuous representative on almost every
line parallel to either coordinate axis. Its derivative vanishes almost
everywhere on its zero set on each such line; Fubini's theorem gives
the asserted equality of the weak gradients on \(G\).

The weak maximal estimate bounds the area of the set where the field
was changed:
\[
 |G^c|\le \frac{C}{B^2}\norm u_2^2.
\]
Combining this with \(w=u\) on \(G\) and the global bound
\(\norm w_\infty\le CB\), we obtain
\[
 \norm w_2^2
 =\int_G|u|^2+\int_{G^c}|w|^2
 \le\norm u_2^2+CB^2|G^c|
 \le C\norm u_2^2.
\]
This calculation also applies when \(G\) is empty. Finally, the
initial mean removal changes the input bounds by at most the fixed
factors already recorded. The resulting \(w\) already has the
divergence of the original fields, so no constant field is added back.
\end{proof}

\begin{remark}
The exceptional set is determined by the field \(u\), whose energy is
small, while the truncation level is determined by \(\norm v_\infty\).
The product of the squared truncation level and the area of this set is
at most \(C\norm u_2^2\), which is why the modification preserves the
energy bound. Modifying a scalar stream function preserves divergence
exactly. No bound on \(\norm v_2\) in terms of \(\norm u_2\), or on
\(\norm u_\infty\), is required.

This is an application of classical Lipschitz truncation.
Schiffer's periodic truncation theorem for closed differential forms
\cite[Theorem~5.1]{Sch22} implies the same simultaneous bounds for
divergence solutions on a torus of any fixed dimension. The proof
above gives an elementary planar argument and makes the roles of
the two input fields explicit.
\end{remark}

\subsection{From a distance cutoff to a bounded flux}
\label{sec:bounded-flux}

We next turn a coupling with uniformly bounded displacement into a
bounded solution of the divergence equation. The natural flux follows
the shortest paths joining the coupled points. These paths can
concentrate at intermediate times, so we first average their flux at
the scale of the distance cutoff. The endpoint density bound controls
the averaged field, and a local correction then restores its exact
divergence.

\begin{lemma}[A bounded divergence solution from a distance cutoff]
\label{lem:bounded-flux}
Let \(\rho_0,\rho_1\) be probability densities on \(\T\) satisfying
\(0\le\rho_0,\rho_1\le M\) almost everywhere. Suppose a coupling
\(\gamma\) of \(\rho_0\,\mathrm dx\) and
\(\rho_1\,\mathrm dx\) is supported on
\(\{(x,y):d_{\T}(x,y)\le r\}\), where \(0<r<1/8\).
There exists \(v\in L^\infty(\T;\R^2)\) such that
\[
 \diver v=\rho_0-\rho_1,\qquad
 \norm v_\infty\le CMr.
\]
The divergence identity holds in the sense of distributions.
\end{lemma}

\begin{proof}
Set \(f=\rho_0-\rho_1\). For \(\gamma\)-almost every \((x,y)\),
let \(\delta(x,y)\in\R^2\) be the shortest displacement vector from
\(x\) to \(y\), so that \(y=x+\delta(x,y)\) modulo \(2\Z^2\).
The restriction on \(r\) makes this vector unique, with
\(\abs{\delta(x,y)}=d_{\T}(x,y)\le r\). All points on the path
\(x+t\delta(x,y)\), \(0\le t\le1\), are interpreted on the torus.

Define the averaged path flux by
\[
 V(z)=\int_{\T\times\T}
 \delta(x,y)\int_0^1
 \eta_r\bigl(z-x-t\delta(x,y)\bigr)\,\mathrm dt
 \,\mathrm d\gamma(x,y).
\]
For fixed \(r\), the kernel and all its derivatives are bounded, so
\(V\) is smooth and differentiation under the integral is valid.
The identity
\[
 \delta(x,y)\cdot\nabla_z
 \eta_r\bigl(z-x-t\delta(x,y)\bigr)
 =-\frac{\mathrm d}{\mathrm dt}
 \eta_r\bigl(z-x-t\delta(x,y)\bigr)
\]
therefore gives
\begin{align*}
 \diver V(z)
 &=\int_{\T\times\T}
   \bigl(\eta_r(z-x)-\eta_r(z-y)\bigr)
   \,\mathrm d\gamma(x,y)\\
 &=(\eta_r*\rho_0)(z)-(\eta_r*\rho_1)(z)
  =(\eta_r*f)(z).
\end{align*}
Here the second equality uses the two marginals of \(\gamma\).

To bound \(V\), fix \(z\). If a kernel in its defining integral is
nonzero, then
\(d_{\T}(z,x+t\delta(x,y))\le r\), and the triangle inequality
implies \(d_{\T}(z,x)\le2r\). Since
\(\norm{\eta_r}_\infty\le Cr^{-2}\), we obtain
\begin{align*}
 |V(z)|
 &\le Cr^{-2}r\,
   \gamma\bigl(B_{\T}(z,2r)\times\T\bigr)\\
 &=Cr^{-1}\int_{B_{\T}(z,2r)}\rho_0(x)\,\mathrm dx
 \le CMr.
\end{align*}
In the last step we used the area bound
\(\abs{B_{\T}(z,2r)}\le Cr^2\). Thus the estimate only involves
the mass of starting endpoints near \(z\); no density bound on the
intermediate path locations is needed.

It remains to correct the difference between \(f\) and its
convolution. Define the periodic vector field
\[
 A_f(z)=\int_{\R^2}\eta_r^{\R}(a)\,a
 \int_0^1 f(z-ta)\,\mathrm dt\,\mathrm da.
\]
This expression defines an essentially bounded field because
\(f\in L^\infty\) and \(\eta_r^{\R}\) is supported on
\(\abs a\le r\). We verify its divergence directly, so that no
derivative of \(f\) is required. For a smooth periodic test function
\(\varphi\), translation invariance of Lebesgue measure and Fubini's
theorem give
\begin{align*}
 -\int_{\T}A_f(z)\cdot\nabla\varphi(z)\,\mathrm dz
 &=-\int_{\R^2}\eta_r^{\R}(a)
   \int_{\T}f(q)\int_0^1
   a\cdot\nabla\varphi(q+ta)\,\mathrm dt\,\mathrm dq\,\mathrm da\\
 &=\int_{\R^2}\eta_r^{\R}(a)
   \int_{\T}f(q)\bigl(\varphi(q)-\varphi(q+a)\bigr)
   \,\mathrm dq\,\mathrm da\\
 &=\int_{\T}\bigl(f(z)-(\eta_r*f)(z)\bigr)
   \varphi(z)\,\mathrm dz.
\end{align*}
For the last equality, translate \(z=q+a\) in the term containing
\(\varphi(q+a)\) and use the definition of periodization. Consequently,
\(\diver A_f=f-\eta_r*f\) in distributions. This identity uses only
the unit integral of the mollifier and is valid without any symmetry
assumption on it.

Finally, \(\norm f_\infty\le M\), and hence
\[
 \norm{A_f}_\infty
 \le\norm f_\infty
     \int_{\R^2}\eta_r^{\R}(a)\abs a\,\mathrm da
 \le Mr.
\]
The field \(v=V+A_f\) has divergence \(f\) and satisfies
\(\norm v_\infty\le CMr\), as claimed.
\end{proof}

\subsection{A single coupling from both flux bounds}
\label{sec:flow}

The final deterministic input converts a vector field with controlled
energy and supremum norm into one coupling with the corresponding
cost and distance bounds. A positive lower bound on the two endpoint
densities permits a linear interpolation between them. Dividing the
fixed flux by this interpolated density gives a velocity field whose
speed and quadratic action are both controlled. This is the
Dacorogna--Moser interpolation used for quadratic transport, as in
\cite[Proposition~2.3]{AST19}; we record its simultaneous
distance-cutoff conclusion and the approximation argument needed here.

\begin{lemma}[A coupling with controlled distance and quadratic cost]
\label{lem:flow}
Let \(\rho_0,\rho_1\) be probability densities on \(\T\) satisfying
\(\rho_0,\rho_1\ge a>0\) almost everywhere. Suppose
\(w\in L^2(\T;\R^2)\cap L^\infty(\T;\R^2)\) satisfies
\(\diver w=\rho_0-\rho_1\) in distributions. Then there is a coupling
\(\Gamma\) of \(\rho_0\,\mathrm dx\) and
\(\rho_1\,\mathrm dx\) such that
\[
 \supp\Gamma\subseteq
 \{(x,y):d_{\T}(x,y)\le a^{-1}\norm w_\infty\},
 \qquad
 \int d_{\T}(x,y)^2\,\mathrm d\Gamma(x,y)
 \le a^{-1}\norm w_2^2.
\]
\end{lemma}

\begin{proof}
First suppose that the densities and the field are smooth. For
\(0\le t\le1\), set
\[
 \rho_t=(1-t)\rho_0+t\rho_1,\qquad
 b_t=\frac{w}{\rho_t}.
\]
The interpolation satisfies \(\rho_t\ge a\) and has total mass one
at every time. Moreover, \(\rho_tb_t=w\), so the divergence assumption
is exactly the continuity equation
\[
 \partial_t\rho_t+\diver(\rho_tb_t)
 =(\rho_1-\rho_0)+\diver w=0.
\]
The field \(b_t\) is smooth on the compact torus, and hence its flow
\(F_t\) exists for \(0\le t\le1\) and consists of diffeomorphisms,
with \(F_0(x)=x\) and
\(\partial_tF_t(x)=b_t(F_t(x))\).

We check that the flow carries the prescribed densities. Its Jacobian
satisfies
\(\partial_t\det DF_t(x)
 =(\diver b_t)(F_t(x))\det DF_t(x)\). Thus the continuity equation
implies
\[
 \frac{\mathrm d}{\mathrm dt}
 \bigl(\rho_t(F_t(x))\det DF_t(x)\bigr)
 =\bigl(\partial_t\rho_t+\diver(\rho_tb_t)\bigr)(F_t(x))
  \det DF_t(x)=0.
\]
Evaluating at \(t=0\) gives
\(\rho_t(F_t(x))\det DF_t(x)=\rho_0(x)\). Change of variables
therefore yields
\((F_t)_\#(\rho_0\,\mathrm dx)=\rho_t\,\mathrm dx\).
In particular,
\[
 \Gamma=(\mathrm{id},F_1)_\#(\rho_0\,\mathrm dx)
\]
is a coupling of the required endpoint measures.

The geodesic distance between the endpoints of a trajectory is at most
its length. Since \(\norm{b_t}_\infty\le a^{-1}\norm w_\infty\),
every initial point satisfies
\[
 d_{\T}(x,F_1(x))
 \le\int_0^1|b_t(F_t(x))|\,\mathrm dt
 \le a^{-1}\norm w_\infty.
\]
This gives the support constraint for \(\Gamma\). For the same
trajectory, Cauchy--Schwarz on the unit time interval gives
\(d_{\T}(x,F_1(x))^2
 \le\int_0^1|b_t(F_t(x))|^2\,\mathrm dt\).
Integrating this inequality against the initial density and using
the pushforward identity at each time yields
\begin{align*}
 \int d_{\T}(x,y)^2\,\mathrm d\Gamma(x,y)
 &\le\int_0^1\int_{\T}
       |b_t(F_t(x))|^2\rho_0(x)\,\mathrm dx\,\mathrm dt\\
 &=\int_0^1\int_{\T}|b_t(y)|^2\rho_t(y)
       \,\mathrm dy\,\mathrm dt\\
 &=\int_0^1\int_{\T}\frac{|w(y)|^2}{\rho_t(y)}
       \,\mathrm dy\,\mathrm dt
 \le a^{-1}\norm w_2^2.
\end{align*}
The support and cost estimates thus hold for the same flow coupling.

For general data, let
\(\rho_i^\eps=\eta_\eps*\rho_i\), \(i=0,1\), and
\(w^\eps=\eta_\eps*w\). Convolving with the same nonnegative
probability kernel preserves total mass, the lower bound
\(\rho_i^\eps\ge a\), and the identity
\(\diver w^\eps=\rho_0^\eps-\rho_1^\eps\).
Young's inequality also gives
\(\norm{w^\eps}_2\le\norm w_2\) and
\(\norm{w^\eps}_\infty\le\norm w_\infty\).
The smooth argument therefore produces couplings \(\Gamma_\eps\)
whose support and quadratic cost satisfy the stated bounds with the
original norms of \(w\), uniformly in \(\eps\).

Because \(\T\times\T\) is compact, a subsequence of
\(\Gamma_\eps\) converges weakly to a probability measure \(\Gamma\)
as \(\eps\downarrow0\). The approximate-identity property gives
\(\rho_i^\eps\to\rho_i\) in \(L^1(\T)\); testing against continuous
functions of either coordinate identifies the marginals of \(\Gamma\)
as \(\rho_0\,\mathrm dx\) and \(\rho_1\,\mathrm dx\).
The common constraint set
\(\{(x,y):d_{\T}(x,y)\le a^{-1}\norm w_\infty\}\) is closed,
so the limit is supported on this set. Finally,
\(d_{\T}(x,y)^2\) is bounded and continuous, and hence its integral
converges along the subsequence. The quadratic-cost bound passes to
\(\Gamma\) as well.
\end{proof}

Combining the three lemmas gives the deterministic statement used in
the proof of the main theorem. The upper density bound converts the
initial distance cutoff into a bounded flux, while the lower density
bound controls the speed and action of the final interpolation.

\begin{proposition}[Combining quadratic energy and bounded displacement]
\label{prop:continuous}
Let \(\rho_0,\rho_1\) be probability densities on \(\T\) with
\(a\le\rho_0,\rho_1\le M\) almost everywhere, where \(a>0\).
Suppose a coupling of these densities is supported on
\(\{(x,y):d_{\T}(x,y)\le r\}\) for some \(0<r<1/8\), and suppose
\(u\in L^2(\T;\R^2)\) satisfies
\(\diver u=\rho_0-\rho_1\) in distributions.
There is a coupling \(\Gamma\) such that
\[
 \supp\Gamma\subseteq
 \{(x,y):d_{\T}(x,y)\le C(M/a)r\},
 \qquad
 \int d_{\T}(x,y)^2\,\mathrm d\Gamma(x,y)
 \le\frac Ca\norm u_2^2.
\]
\end{proposition}

\begin{proof}
Lemma~\ref{lem:bounded-flux} supplies a field \(v\) with
\(\diver v=\rho_0-\rho_1\) and \(\norm v_\infty\le CMr\).
Apply Lemma~\ref{lem:interpolate} to \(u\) and \(v\) to obtain a
field \(w\) satisfying
\[
 \diver w=\rho_0-\rho_1,\qquad
 \norm w_2^2\le C\norm u_2^2,\qquad
 \norm w_\infty\le CMr.
\]
Lemma~\ref{lem:flow}, applied to this single field, gives a coupling
whose distance cutoff is at most \(C(M/a)r\) and whose quadratic
cost is at most \(Ca^{-1}\norm u_2^2\).
\end{proof}

\section{Regularizing independent uniform samples}
\label{sec:random-inputs}

We now verify the hypotheses of Proposition~\ref{prop:continuous} for
smoothed empirical measures of the torus samples in
Theorem~\ref{thm:torus}. We use \(\Pbb,\Ebb\) for their joint law
and expectation. With the mollifier fixed in
\S\ref{sec:preliminaries}, set
\[
 \mu_X=\frac1N\sum_{i=1}^N\delta_{X_i},\qquad
 \mu_Y=\frac1N\sum_{j=1}^N\delta_{Y_j},\qquad
 \rho_0=\eta_h*\mu_X,\qquad
 \rho_1=\eta_h*\mu_Y,
 \qquad
 h=A\sqrt{\frac{\log N}{N}}.
\]
We choose \(A\) once and for all as in Lemma~\ref{lem:density}, and
take \(N\) sufficiently large that \(h<1/8\). The three estimates
below supply, respectively, the positive density bounds, a divergence
solution with small energy, and a coupling with a distance cutoff.

\subsection{Uniform bounds on the densities}

\begin{lemma}
\label{lem:density}
For a sufficiently large fixed \(A\),
\[
 \Pbb\left(\frac18\le\rho_0,\rho_1\le\frac38
 \text{ everywhere on }\T\right)\ge1-CN^{-2}
\]
for all sufficiently large \(N\).
\end{lemma}

\begin{proof}
For a fixed \(z\), the summands \(\eta_h(z-X_i)\) are independent,
with
\[
 \Ebb\eta_h(z-X_i)=\frac14,\qquad
 0\le\eta_h(z-X_i)\le Ch^{-2},\qquad
 \Var(\eta_h(z-X_i))\le Ch^{-2}.
\]
For example, the variance bound follows from
\(\Ebb\eta_h(z-X_i)^2=\frac14\int_{\T}\eta_h^2\le Ch^{-2}\).
Bernstein's inequality gives
\[
 \Pbb\left(\left|\rho_0(z)-\frac14\right|>\frac1{16}\right)
 \le 2\exp(-cNh^2).
\]

Deterministically, \(\norm{\nabla\rho_0}_\infty\le Ch^{-3}\).
Choose a torus net of mesh \(c_1h^3\), with \(c_1>0\) sufficiently
small. It has at most \(Ch^{-6}\) points, and the density changes
by at most \(1/16\) between an arbitrary point and its nearest
net point. A union bound therefore gives
\[
 \Pbb\left(\sup_{z\in\T}
 \left|\rho_0(z)-\frac14\right|>\frac18\right)
 \le Ch^{-6}\exp(-cNh^2).
\]
The same argument applies to \(\rho_1\). Since
\(Nh^2=A^2\log N\) and \(h^{-6}\le C_A N^3\), choosing \(A\)
large enough makes the combined failure probability at most
\(CN^{-2}\).
\end{proof}

\subsection{A quadratic-energy estimate}

\begin{lemma}
\label{lem:energy}
Let \(f=\rho_0-\rho_1\), and define
\[
 u=-\nabla(-\Delta)^{-1}f
\]
using the inverse Laplacian on mean-zero periodic functions. Then
\(\overline u=0\), \(\diver u=f\), and, for a fixed sufficiently
large \(C\),
\[
 \Pbb\left(\norm u_2^2>C\frac{\log N}{N}\right)
 \le \frac{C}{(\log N)^2}.
\]
\end{lemma}

\begin{proof}
The inverse Laplacian is well defined because both densities have
integral one, so \(f\) has zero mean. We estimate both the mean and
the variance of the energy. With the
Fourier convention fixed in \S\ref{sec:interpolation}, Parseval's
identity reads \(\norm g_2^2=4\sum_{k\in\Z^2}|\widehat g(k)|^2\).
For the fixed function \(\eta\) on \(\R^2\), write
\(\widehat\eta(\xi)=\int_{\R^2}\eta(z)e^{-i\xi\cdot z}\,\mathrm dz\).
The Fourier coefficient of its periodic smoothing kernel is
\(\widehat{\eta_h}(k)=\frac14\widehat\eta(\pi hk)\).

Pool the \(2N\) points as \(Z_1,\ldots,Z_{2N}\), and put
\(\eps_a=1\) for \(a\le N\) and \(\eps_a=-1\) for \(a>N\).
Then
\[
 \widehat f(k)=
 \frac{\widehat\eta(\pi hk)}{4N}
 \sum_{a=1}^{2N}\eps_a e^{-i\pi k\cdot Z_a},
 \qquad
 \widehat f(0)=0.
\]
The eigenvalue of \(-\Delta\) at nonzero frequency \(k\) is
\(\pi^2|k|^2\). Parseval's identity therefore gives
\begin{equation}
\label{eq:energy-kernel}
 E_h:=\norm u_2^2
 =\frac1{N^2}\sum_{a,b=1}^{2N}
   \eps_a\eps_b K_h(Z_a-Z_b),
\end{equation}
where the real, even, mean-zero kernel \(K_h\) is
\[
 K_h(z)=\frac1{4\pi^2}
 \sum_{k\in\Z^2\setminus\{0\}}
 \frac{|\widehat\eta(\pi hk)|^2}{|k|^2}e^{i\pi k\cdot z}.
\]

The coefficients of \(K_h\) are unchanged by \(k\mapsto-k\), since
\(\eta\) is real and
\(\widehat\eta(-\xi)=\overline{\widehat\eta(\xi)}\). Thus \(K_h\)
is real and even without requiring \(\eta\) itself to be symmetric.
Since \(|\widehat\eta|\le1\) and \(\widehat\eta\) decays faster than
any power, the two-dimensional lattice sum gives
\[
 K_h(0)\le C\bigl(1+\log(1/h)\bigr).
\]
Indeed, summing \(|k|^{-2}\) over \(1\le|k|\le h^{-1}\)
gives \(O(1+\log(1/h))\). Beyond \(h^{-1}\), the Fourier
decay bounds the remaining sum by a constant.

The diagonal part of \eqref{eq:energy-kernel} is the deterministic
quantity \(2K_h(0)/N\). All off-diagonal terms have expectation
zero, because \(Z_a-Z_b\) is uniform on \(\T\) and \(K_h\) has
zero mean. Hence
\[
 \Ebb E_h=\frac{2K_h(0)}N
 \le C\frac{\log N}{N}.
\]

An expectation estimate alone would not yield the desired conclusion
with a fixed constant. To control its fluctuations, let \(Z,Z'\)
be independent uniform torus points. Parseval gives
\[
 \Ebb K_h(Z-Z')^2
 =\frac1{16\pi^4}
 \sum_{k\ne0}\frac{|\widehat\eta(\pi hk)|^4}{|k|^4}
 \le C.
\]
Distinct unordered pairs \(\{a,b\}\) and \(\{c,d\}\) give
uncorrelated off-diagonal terms. They are independent if they
have no common index. If, for example, they share \(a=c\)
and their other indices differ, conditioning on \(Z_a\) shows
that both conditional means are zero and the remaining variables
are independent. Since \(K_h\) is even, we obtain exactly
\[
 \Var(E_h)
 =\frac4{N^4}\binom{2N}{2}\Ebb K_h(Z-Z')^2
 \le \frac C{N^2}.
\]
Choose the threshold constant strictly larger than the constant in
the expectation bound. Chebyshev's inequality then gives
\[
 \Pbb\left(E_h>C\frac{\log N}{N}\right)
 \le\frac{C/N^2}{(\log N/N)^2}
 \le\frac C{(\log N)^2},
\]
as required.
\end{proof}

\subsection{The bottleneck input and common smoothing}

\begin{lemma}[Bottleneck matching input]
\label{lem:bottleneck}
There are universal constants \(b,c,C>0\) such that independent
uniform samples on \(\T\) satisfy
\[
 \Pbb\left(
 \exists \sigma\in\mathfrak S_N:
 \max_i d_{\T}(X_i,Y_{\sigma(i)})
 \le b\,\frac{(\log N)^{3/4}}{\sqrt N}
 \right)
 \ge1-C\exp\bigl(-c(\log N)^{3/2}\bigr).
\]
\end{lemma}

\begin{proof}
We first work in the unit square. The Leighton--Shor theorem
\cite{LS89}, as stated in \cite[Theorem~4.7.1]{Tal21}, matches an
independent uniform sample to any fixed evenly spread family of \(N\)
points with radius at most \(C(\log N)^{3/4}/\sqrt N\) and failure
probability at most \(C\exp(-c(\log N)^{3/2})\).

We specify a reference family valid for every \(N\). In the definition
of \cite[\S4.3]{Tal21}, the square is partitioned into \(N\) rectangles
of area \(1/N\), each containing its prescribed point and lying within
distance \(10/\sqrt N\) of that point. To construct such a partition,
let \(m=\lfloor\sqrt N\rfloor\) and choose integers
\(n_1,\ldots,n_m\), differing by at most one, whose sum is \(N\).
Make row \(j\) have height \(n_j/N\), and divide it into \(n_j\)
rectangles of width \(1/n_j\). Each rectangle has area \(1/N\).
Since \(m\le n_j\le m+2\), its width is at most \(2/\sqrt N\)
and its height is at most \(3/\sqrt N\). Its diameter is therefore
at most \(\sqrt{13}/\sqrt N\), so the rectangle centers meet the
cited definition.

Apply the theorem to each sample using these same centers. If
\(\tau_X,\tau_Y\) map the two samples to the reference family, then
\(\sigma=\tau_Y^{-1}\circ\tau_X\) matches the samples to each other.
The triangle inequality increases the radius by at most a factor two,
and a union bound adds the two failure probabilities.

Finally, represent the uniform torus points in \([-1,1)^2\) and map
this square to \([0,1)^2\) by scaling and translation. The preceding
argument gives a matching whose Euclidean edge lengths in the original
representatives are bounded by twice the unit-square bound. Torus
distance is at most the Euclidean distance between representatives.
Absorbing these fixed factors into \(b,c,C\) proves the lemma.
\end{proof}

For the remainder of the proof, let
\[
 r=b\,\frac{(\log N)^{3/4}}{\sqrt N}.
\]
For all sufficiently large \(N\), we have \(r<1/8\).
If \(\sigma\) is the matching in Lemma~\ref{lem:bottleneck},
define a coupling of the smoothed densities by using the same
shift at its two endpoints:
\begin{equation}
\label{eq:common-shift}
 \gamma=\frac1N\sum_{i=1}^N
 \int_{\R^2}\eta_h^{\R}(z)\,
 \delta_{(X_i+z,\;Y_{\sigma(i)}+z)}\,\mathrm dz.
\end{equation}
All points are understood modulo \(2\Z^2\).
The marginals are \(\rho_0\,\mathrm dx,\rho_1\,\mathrm dx\).
Translation invariance gives
\(d_{\T}(X_i+z,Y_{\sigma(i)}+z)=d_{\T}(X_i,Y_{\sigma(i)})\),
so \(\gamma\) has the same cutoff \(r\).

\section{Returning to a permutation}
\label{sec:rounding}

To complete the proof of Theorem~\ref{thm:torus}, we pass from a
coupling of the smoothed densities to a matching of the original torus
sample points. The conditional weights below assign each
smoothed location to its empirical source. This gives a fractional
matching supported on short edges. The Birkhoff--von Neumann theorem
then selects a permutation while preserving that support and controlling
the quadratic cost by averaging.

\begin{lemma}[Rounding with a distance cutoff]
\label{lem:rounding}
Let \(X_i,Y_j\in\T\), and let
\(\rho_0=\eta_h*\mu_X\) and \(\rho_1=\eta_h*\mu_Y\) be positive.
Suppose a coupling \(\Gamma\) of these densities satisfies
\[
 \supp\Gamma\subseteq\{d_{\T}(x,y)\le D\},
 \qquad
 \int d_{\T}(x,y)^2\,\mathrm d\Gamma(x,y)\le E.
\]
There is a permutation \(\pi\in\mathfrak S_N\) such that
\[
 \max_i d_{\T}(X_i,Y_{\pi(i)})\le2h+D,\qquad
 \frac1N\sum_i d_{\T}(X_i,Y_{\pi(i)})^2
 \le3(2h^2+E).
\]
\end{lemma}

\begin{proof}
The conditional weights
\[
 \alpha_i(x)=\frac{\eta_h(x-X_i)}{N\rho_0(x)},
 \qquad
 \beta_j(y)=\frac{\eta_h(y-Y_j)}{N\rho_1(y)}
\]
are nonnegative and satisfy
\(\sum_i\alpha_i(x)=\sum_j\beta_j(y)=1\).
Define the matrix
\[
 p_{ij}=\int_{\T\times\T}
 \alpha_i(x)\beta_j(y)\,\mathrm d\Gamma(x,y).
\]
Its row and column sums are \(1/N\). For example,
\[
 \sum_jp_{ij}=\int_{\T}\alpha_i(x)\rho_0(x)\,\mathrm dx
 =\frac1N\int_{\T}\eta_h(x-X_i)\,\mathrm dx=\frac1N.
\]

Wherever \(\alpha_i(x)\beta_j(y)>0\), the smoothing kernel implies
\(d_{\T}(X_i,x)\le h\) and \(d_{\T}(Y_j,y)\le h\). Thus
\[
 p_{ij}>0\quad\Longrightarrow\quad
 d_{\T}(X_i,Y_j)\le2h+D.
\]
The squared triangle inequality, integrated against the same
conditional weights, gives
\begin{align*}
 \sum_{i,j}p_{ij}d_{\T}(X_i,Y_j)^2
 &\le3\int\sum_{i,j}\alpha_i(x)\beta_j(y)
 \bigl(d_{\T}(X_i,x)^2+d_{\T}(x,y)^2
       +d_{\T}(y,Y_j)^2\bigr)\,\mathrm d\Gamma\\
 &\le3(2h^2+E).
\end{align*}

The matrix \(Np\) is doubly stochastic. By the
Birkhoff--von Neumann theorem \cite{Bir46}, it is a convex combination
of permutation matrices. Every permutation with positive weight uses
only edges for which \(p_{ij}>0\): a zero matrix entry cannot be a
positive convex combination of nonnegative entries. Therefore all
permutations in the decomposition satisfy the stated cutoff.
At least one of them has normalized quadratic cost at most
\(\sum_{i,j}p_{ij}d_{\T}(X_i,Y_j)^2\), because that sum is the
weighted average of their normalized costs.
\end{proof}

\begin{proof}[Proof of Theorem~\ref{thm:torus}]
Intersect the events of Lemmas~\ref{lem:density},
\ref{lem:energy}, and \ref{lem:bottleneck}. Their intersection has
probability at least
\[
 1-CN^{-2}-\frac C{(\log N)^2}
   -C\exp\bigl(-c(\log N)^{3/2}\bigr)
 \ge1-\frac C{(\log N)^2}
\]
for all sufficiently large \(N\). No independence between these
three events is required.

On this event, \(\rho_0,\rho_1\) lie between \(a=1/8\) and
\(M=3/8\), the field \(u\) has energy at most \(C\log N/N\),
and \eqref{eq:common-shift} is a coupling with cutoff \(r\).
Proposition~\ref{prop:continuous} gives a single coupling
\(\Gamma\) such that
\[
 \supp\Gamma\subseteq\{d_{\T}(x,y)\le Cr\},
 \qquad
 \int d_{\T}(x,y)^2\,\mathrm d\Gamma(x,y)
 \le C\frac{\log N}{N}.
\]
Apply Lemma~\ref{lem:rounding}. Since
\[
 h^2=A^2\frac{\log N}{N},
 \qquad
 \frac hr=\frac{A}{b(\log N)^{1/4}}\longrightarrow0,
\]
it gives one permutation satisfying
\[
 \sum_i d_{\T}(X_i,Y_{\pi(i)})^2\le C\log N,
 \qquad
 \max_i d_{\T}(X_i,Y_{\pi(i)})\le Cr.
\]
Since \(r=b(\log N)^{3/4}/\sqrt N\), choosing fixed
\(C_0,C_1,C_2\) large enough in the cost, cutoff, and
failure-probability bounds proves Theorem~\ref{thm:torus}.
\end{proof}

\section{Expected cost and evenly spread targets}
\label{sec:further}

The preceding argument also controls the expected cost of a single
matching while retaining the stronger probability estimate supplied by
the bottleneck theorem. The same method applies when one sample is
replaced by a fixed evenly spread family. In both cases, a finite
optimization rule makes the choice of permutation measurable.

\subsection{Expected cost and the bottleneck probability}

\begin{corollary}
\label{cor:expected}
In the square model of Theorem~\ref{thm:main}, there are universal
constants \(C,c>0\) and a measurable choice of
\(\pi\in\mathfrak S_N\) such that, for all sufficiently large \(N\),
\[
 \Ebb\sum_{i=1}^N d(X_i,Y_{\pi(i)})^2\le C\log N,
\]
and
\[
 \Pbb\left(
  \max_i d(X_i,Y_{\pi(i)})>
  C\frac{(\log N)^{3/4}}{\sqrt N}
 \right)
 \le C\exp\bigl(-c(\log N)^{3/2}\bigr).
\]
The same choice of \(\pi\) satisfies both bounds in
Theorem~\ref{thm:main} with probability at least
\(1-C/(\log N)^2\). In particular,
\[
 \Ebb\sum_{i=1}^N d(X_i,Y_{\pi(i)})
 \le C\sqrt{N\log N}.
\]
\end{corollary}

\begin{proof}
Recall \(r=b(\log N)^{3/4}/\sqrt N\), and fix a sufficiently large
universal constant \(L\). For each pair of square samples, choose a
permutation minimizing the total squared cost among those whose longest
edge is at most \(Lr\). Break ties by a fixed lexicographic order, and
choose the identity permutation if this constraint is infeasible.
There are only finitely many permutations, and the feasibility and cost
of each are continuous inequalities in the sample coordinates.
Consequently, this rule is measurable. We choose \(L\) below, uniformly
in \(N\).

To analyze the rule, let
\(Z_1,\ldots,Z_N,Z'_1,\ldots,Z'_N\) be independent uniform torus
points. The arrays \((F(Z_i))_{i=1}^N\) and
\((F(Z'_i))_{i=1}^N\) have the joint law of the square samples.
Let \(\rho_0,\rho_1\) and \(E_h\) denote their torus smoothing
densities and Poisson energy, as in \S\ref{sec:random-inputs}. Let
\(\mathcal B\) be the torus bottleneck event from
Lemma~\ref{lem:bottleneck}, and set
\[
 \mathcal D=\left\{\frac18\le\rho_0,\rho_1\le\frac38\right\}.
\]
The existing estimates give
\[
 \Pbb(\mathcal D^c)\le CN^{-2},\qquad
 \Pbb(\mathcal B^c)
 \le C\exp\bigl(-c(\log N)^{3/2}\bigr),\qquad
 \Ebb E_h\le C\frac{\log N}{N}.
\]

On \(\mathcal B\cap\mathcal D\), the proof of
Theorem~\ref{thm:torus}, without imposing the energy event, constructs
a torus permutation with cutoff \(Cr+2h\) and total squared cost
at most \(CN(E_h+h^2)\). Indeed,
Proposition~\ref{prop:continuous} bounds the coupling cost by \(CE_h\),
and Lemma~\ref{lem:rounding} then gives this estimate. Folding does
not increase any edge length. Since \(h=o(r)\), a fixed sufficiently
large \(L\) makes this square permutation feasible for the selection
rule, whose cost is therefore at most \(CN(E_h+h^2)\).

On \(\mathcal B\cap\mathcal D^c\), fold the original bottleneck
permutation instead. It has cutoff \(r\) and total squared cost at
most \(Nr^2\), so the selected permutation has no greater cost.
Thus the cutoff constraint is feasible on all of \(\mathcal B\).
On \(\mathcal B^c\), the selected permutation still has total squared
cost at most \(2N\), since the squared diameter of the unit square
is \(2\). Taking expectations gives
\begin{align*}
 \Ebb\sum_i d(X_i,Y_{\pi(i)})^2
 &\le CN(\Ebb E_h+h^2)
      +Nr^2\Pbb(\mathcal D^c)+2N\Pbb(\mathcal B^c)\\
 &\le C\log N.
\end{align*}
Here \(Nh^2=O(\log N)\), \(Nr^2=O((\log N)^{3/2})\), and
\(N\exp(-c(\log N)^{3/2})\to0\).

Feasibility on \(\mathcal B\) proves the claimed bottleneck
probability. On
\(\mathcal B\cap\mathcal D\cap\{E_h\le C\log N/N\}\), the
same selected permutation also has total squared cost at most
\(C\log N\). Lemmas~\ref{lem:density}, \ref{lem:energy}, and
\ref{lem:bottleneck} bound the probability of its complement by
\(C/(\log N)^2\). Finally, Cauchy--Schwarz, first for the sum and
then for expectation, gives
\[
 \Ebb\sum_i d(X_i,Y_{\pi(i)})
 \le\sqrt{N\,\Ebb\sum_i d(X_i,Y_{\pi(i)})^2}
 \le C\sqrt{N\log N}.
\]
\end{proof}

\subsection{Matching to fixed evenly spread points}

We call a fixed family \(Y_1,\ldots,Y_N\in[0,1]^2\)
\emph{evenly spread} if there is a measurable partition
\(Q_1,\ldots,Q_N\) of the square such that
\[
 |Q_j|=\frac1N,\qquad
 \sup_{y\in Q_j}|y-Y_j|\le\frac{10}{\sqrt N},
 \qquad 1\le j\le N.
\]
This includes the rectangular partitions in the definition used in
\cite[\S4.3]{Tal21}. Only the displayed area and distance conditions
are needed below.

\begin{corollary}
\label{cor:grid}
Let \(Y_1,\ldots,Y_N\) be a fixed evenly spread family, and let
\(X_1,\ldots,X_N\) be independent uniform points in the unit square.
There are universal constants \(C,c>0\), independent of the fixed
family, and a measurable choice \(\pi=\pi(X_1,\ldots,X_N)\) such
that, for all sufficiently large \(N\),
\[
 \Pbb\left(
  \sum_i d(X_i,Y_{\pi(i)})^2\le C\log N,
  \quad
  \max_i d(X_i,Y_{\pi(i)})
  \le C\frac{(\log N)^{3/4}}{\sqrt N}
 \right)
 \ge1-\frac{C}{(\log N)^2}.
\]
The same choice satisfies
\[
 \Ebb\sum_i d(X_i,Y_{\pi(i)})^2\le C\log N,
 \qquad
 \Ebb\sum_i d(X_i,Y_{\pi(i)})\le C\sqrt{N\log N},
\]
and
\[
 \Pbb\left(
  \max_i d(X_i,Y_{\pi(i)})>
  C\frac{(\log N)^{3/4}}{\sqrt N}
 \right)
 \le C\exp\bigl(-c(\log N)^{3/2}\bigr).
\]
\end{corollary}

\begin{proof}
Let \(Z_1,\ldots,Z_N\) be independent uniform torus points, so that
\((F(Z_i))_{i=1}^N\) has the law of the square sample. We first
couple the smoothed torus sample to uniform torus measure. Set
\[
 \rho_0=\eta_h*\left(\frac1N\sum_i\delta_{Z_i}\right),
 \qquad \rho_1=\frac14,\qquad
 u=-\nabla(-\Delta)^{-1}(\rho_0-\tfrac14),
 \qquad E_h=\norm u_2^2.
\]
The proof of Lemma~\ref{lem:density} gives
\[
 \Pbb(\mathcal D^c)\le CN^{-2},
 \qquad
 \mathcal D=\left\{\frac18\le\rho_0\le\frac38\right\}.
\]
The Fourier calculation in Lemma~\ref{lem:energy}, with \(N\)
unsigned points in place of \(2N\) signed points, gives
\[
 E_h=\frac1{N^2}\sum_{a,b=1}^N K_h(Z_a-Z_b).
\]
The kernel \(K_h\) is the same real, even, mean-zero kernel as
before. Its diagonal contribution is deterministic, and distinct
unordered off-diagonal pairs are uncorrelated. Consequently,
\[
 \Ebb E_h=\frac{K_h(0)}N\le C\frac{\log N}{N},
 \qquad
 \Var(E_h)=\frac4{N^4}\binom N2
       \Ebb K_h(Z-Z')^2\le\frac C{N^2},
\]
where \(Z,Z'\) are independent uniform torus points. Chebyshev's
inequality therefore gives
\[
 \Pbb\left(E_h>C\frac{\log N}{N}\right)
 \le\frac C{(\log N)^2}
\]
for a sufficiently large fixed \(C\).

Write \(s=10/\sqrt N\). The rectangular construction in the proof
of Lemma~\ref{lem:bottleneck}, scaled to the torus, gives a fixed
partition \(\widetilde Q_1,\ldots,\widetilde Q_N\) into cells of
ordinary area \(4/N\), with centers \(q_j\), such that
\(d_{\T}(y,q_j)\le Cs\) for \(y\in\widetilde Q_j\).
Apply the one-sample Leighton--Shor theorem
\cite[Theorem~4.7.1]{Tal21} to these centers as in that proof.
After increasing the universal constant \(b\) in the definition of
\(r\), if necessary, this gives an event \(\mathcal B\) with
\[
 \Pbb(\mathcal B^c)
 \le C\exp\bigl(-c(\log N)^{3/2}\bigr)
\]
on which some \(\sigma\in\mathfrak S_N\) satisfies
\(\max_i d_{\T}(Z_i,q_{\sigma(i)})\le r\).
On this event, the measure
\[
 \Lambda=\frac1N\sum_i\delta_{Z_i}\otimes
 \left(\frac N4\mathbf1_{\widetilde Q_{\sigma(i)}}(y)
       \,\mathrm dy\right)
\]
couples the torus empirical measure to \(\mathrm dy/4\) and has
cutoff \(r+Cs\). Average its common translates
\((x,y)\mapsto(x+a,y+a)\) against
\(\eta_h^{\R}(a)\,\mathrm da\). The first marginal becomes
\(\rho_0\,\mathrm dx\), the second remains \(\mathrm dy/4\),
and translation invariance preserves the cutoff \(r+Cs\).

Since \(s=o(r)\) and \(r+Cs<1/8\) for sufficiently large \(N\),
Proposition~\ref{prop:continuous} applies on
\(\mathcal B\cap\mathcal D\). It supplies a coupling \(\Gamma\)
of \(\rho_0\,\mathrm dx\) and \(\mathrm dy/4\) such that
\[
 \supp\Gamma\subseteq\{d_{\T}(x,y)\le Cr\},
 \qquad
 \int d_{\T}(x,y)^2\,\mathrm d\Gamma(x,y)\le CE_h.
\]
To round this coupling to the prescribed square targets, use their
partition \(Q_1,\ldots,Q_N\) and define
\[
 \alpha_i(x)=\frac{\eta_h(x-Z_i)}{N\rho_0(x)},
 \qquad
 p_{ij}=\int\alpha_i(x)\mathbf1_{Q_j}(F(y))
              \,\mathrm d\Gamma(x,y).
\]
The row sums are \(1/N\), as in Lemma~\ref{lem:rounding}. The
column sums are also \(1/N\), because
\[
 \sum_i p_{ij}
 =\frac14\int_{\T}\mathbf1_{Q_j}(F(y))\,\mathrm dy
 =|Q_j|=\frac1N.
\]
Here we used that folding sends uniform torus measure to uniform
square measure. On the support of a contributing weight,
\[
 |F(Z_i)-Y_j|
 \le d_{\T}(Z_i,x)+d_{\T}(x,y)+s
 \le h+Cr+s.
\]
The same squared triangle inequality as before gives
\[
 \sum_{i,j}p_{ij}|F(Z_i)-Y_j|^2
 \le3(h^2+CE_h+s^2).
\]
The Birkhoff argument in Lemma~\ref{lem:rounding} therefore produces
a square permutation with cutoff \(h+Cr+s\) and total squared
cost at most \(CN(E_h+h^2+s^2)\).

We also need a fallback that does not require \(\mathcal D\).
On all of \(\mathcal B\), round \(\Lambda\) directly by using
instead the matrix
\[
 p_{ij}=\frac14\int_{\widetilde Q_{\sigma(i)}}
                   \mathbf1_{Q_j}(F(y))\,\mathrm dy.
\]
Its row and column sums are again \(1/N\). If \(p_{ij}>0\), then
\[
 |F(Z_i)-Y_j|\le r+Cs+s\le Cr.
\]
Thus Birkhoff rounding gives a permutation with cutoff \(Cr\) and
total squared cost at most \(CNr^2\), regardless of the smoothed
density bounds.

Choose the square permutation by minimizing total squared cost under
the cutoff \(Lr\), for a fixed universal \(L\) large enough to
include both constructions above. Break ties lexicographically and
use the identity permutation when the constraint is infeasible. This
is the finite measurable selection rule used in
Corollary~\ref{cor:expected}; in particular, it depends only on
the square sample and the fixed targets, not on a choice of torus
lift. The folded sample law allows us to analyze it using the torus
construction.

On \(\mathcal B\cap\mathcal D\), optimality bounds its cost by
\(CN(E_h+h^2+s^2)\). On
\(\mathcal B\cap\mathcal D^c\), the fallback bounds its cost by
\(CNr^2\); on \(\mathcal B^c\), its cost is at most \(2N\).
It follows that
\begin{align*}
 \Ebb\sum_i d(X_i,Y_{\pi(i)})^2
 &\le CN(\Ebb E_h+h^2+s^2)
       +CNr^2\Pbb(\mathcal D^c)+2N\Pbb(\mathcal B^c)\\
 &\le C\log N.
\end{align*}
The expected total-distance bound follows by Cauchy--Schwarz.
The cutoff constraint is feasible on \(\mathcal B\), which gives
the stated bottleneck probability. Finally, on
\(\mathcal B\cap\mathcal D\cap\{E_h\le C\log N/N\}\), the
same permutation has total squared cost at most \(C\log N\).
The probability of the complement is at most \(C/(\log N)^2\),
proving the remaining assertion. All constants used the fixed target
family only through the bound \(s=10/\sqrt N\), so they are
uniform over all such families.
\end{proof}

\begin{thebibliography}{9}

\bibitem{AST19}
L.~Ambrosio, F.~Stra, and D.~Trevisan,
\emph{A PDE approach to a 2-dimensional matching problem},
Probability Theory and Related Fields \textbf{173} (2019), 433--477.
\href{https://doi.org/10.1007/s00440-018-0837-x}{doi:10.1007/s00440-018-0837-x}.

\bibitem{LS89}
T.~Leighton and P.~Shor,
\emph{Tight bounds for minimax grid matching with applications to the
average-case analysis of algorithms},
Combinatorica \textbf{9} (1989), 161--187.

\bibitem{Tal21}
M.~Talagrand,
\emph{Upper and Lower Bounds for Stochastic Processes:
Decomposition Theorems}, second edition,
Ergebnisse der Mathematik und ihrer Grenzgebiete, vol.~60,
Springer, 2021.
\href{https://michel.talagrand.net/ULBSPRINGER.pdf}{Author's PDF}.

\bibitem{Tal26}
M.~Talagrand,
\emph{My favorite problems}, preprint, 2026.
\href{https://michel.talagrand.net/Korea15.pdf}{Author's PDF},
Problem~5.3 and Conjecture~5.8 in the version cited.

\bibitem{Oui26}
F.~Ouimet,
\emph{A solution to Problem 5.3 in Talagrand's My favorite problems:
Simultaneous sum and bottleneck bounds for random planar matchings},
preprint, 2026, version~2.
\href{https://doi.org/10.13140/RG.2.2.26527.93606/2}{doi:10.13140/RG.2.2.26527.93606/2}.

\bibitem{Sch22}
S.~Schiffer,
\emph{\(L^\infty\)-truncation of closed differential forms},
Calculus of Variations and Partial Differential Equations
\textbf{61} (2022), article~135.
\href{https://doi.org/10.1007/s00526-022-02236-1}{doi:10.1007/s00526-022-02236-1}.

\bibitem{Haj00}
P.~Haj{\l}asz,
\emph{Sobolev mappings, co-area formula and related topics},
in \emph{Proceedings on Analysis and Geometry},
Sobolev Institute Press, Novosibirsk, 2000, pp.~227--254.
\href{https://sites.pitt.edu/~hajlasz/OriginalPublications/Hajlasz-Sobolev-ProceedingsOnAnalysisAndGeometry-2000-227-254.pdf}{Author's PDF}.

\bibitem{Bir46}
G.~Birkhoff,
\emph{Tres observaciones sobre el \'algebra lineal},
Universidad Nacional de Tucum\'an, Revista, Serie A
\textbf{5} (1946), 147--151.

\end{thebibliography}
\end{document}
